\chapter{INTRODUÇÃO E JUSTIFICATIVA}

A compreensão do espaço geográfico contemporâneo exige uma análise que transcenda a dimensão puramente física do território. Segundo \cite{milton-ne}, o espaço é uma construção social, composto por sistemas de objetos e sistemas de ações, que refletem as desigualdades e as relações de poder da sociedade capitalista.

\section*{Justificativa}
Esta pesquisa justifica-se pela necessidade de repensar as práticas pedagógicas no ensino de Geografia. Como aponta \cite{tec-ensi-pqn}, as técnicas de ensino não podem ser utilizadas de forma mecânica; elas precisam dialogar com a realidade dos estudantes e com os desafios do mundo contemporâneo, formando cidadãos críticos e atuantes.